% ============================================================
%  CHAPTER 6 — RESULTS AND DISCUSSION
%  Fully rewritten 2026-08-28 -- see VERIFICATION_NOTES.md.
%  Every number in this chapter is sourced directly from CLAUDE.md's
%  verified, dated result log. The original draft's entire results
%  chapter -- including a full "achieved" column with checkmarks against
%  fabricated targets, and a complete pilot-study burden-comparison table
%  for a study that has never run -- has been replaced, not edited.
% ============================================================
\chapter{Results and Discussion}
\label{chap:results}

This chapter reports this project's real, measured results. Where a result
is still pending or blocked, it is stated as such rather than filled with a
plausible-looking number.

% ------------------------------------------------------------
\section{ISCO-08 Classification: The WISCO External Benchmark}
\label{sec:res-isco}
% ------------------------------------------------------------

\subsection{Canonical Baseline: Flat vs.\ Hierarchical Retrieval}

\begin{table}[htbp]
\centering
\caption{ISCO-08 exact 4-digit accuracy, full 18,747-case WISCO heldout split.}
\label{tab:ch6-wisco-canonical}
\small\renewcommand{\arraystretch}{1.4}
\begin{tabular}{lcc}
\toprule
\textbf{Configuration} & \textbf{Exact 4-digit accuracy} & \textbf{95\% Wilson CI} \\
\midrule
Flat retrieval & \textbf{21.19\%} & [20.61\%, 21.78\%] \\
Hierarchical retrieval & 10.35\% & [9.93\%, 10.80\%] \\
\bottomrule
\end{tabular}
\end{table}

\textbf{Hierarchical retrieval is 10.84 percentage points \emph{worse} than
flat retrieval} (McNemar exact $p \approx 1.86\times10^{-301}$) --- a real,
statistically decisive negative finding. This directly contradicts an
earlier internal draft of this thesis, which reported the opposite
(hierarchical outperforming flat by 6.4pp) without this benchmark having
been run. Hierarchical retrieval is never described as outperforming flat
retrieval anywhere else in this thesis.

\subsection{Improving the Best Configuration: Embedding Model and Catalogue
Enrichment}

Starting from flat retrieval (the better baseline), two independent
improvements were tested on the identical 18,747-case set.

\begin{table}[htbp]
\centering
\caption{Additive improvements to flat retrieval, full heldout set.}
\label{tab:ch6-additive}
\small\renewcommand{\arraystretch}{1.4}
\begin{tabular}{lccc}
\toprule
\textbf{Configuration} & \textbf{Accuracy} & \textbf{$\Delta$ vs.\ baseline} & \textbf{McNemar $p$} \\
\midrule
Baseline (e5-small) & 21.19\% & --- & --- \\
+ e5-large embeddings & 29.70\% & $+8.50$pp & $7.34\times10^{-167}$ \\
+ Official-text enrichment (e5-small) & 32.55\% & $+11.36$pp & $2.19\times10^{-274}$ \\
+ Both combined & \textbf{40.95\%} & \textbf{$+19.75$pp} & $\approx 0$ \\
\bottomrule
\end{tabular}
\end{table}

\textbf{40.95\% is this thesis's headline, best-tested ISCO-08 result} ---
nearly double the 21.19\% original baseline, and the number this thesis
cites wherever an ISCO-08 accuracy figure is needed, not the fabricated
85.3\% an earlier internal draft stated. The two improvements are genuinely
additive (not redundant), consistent with them working through independent
mechanisms: embedding-model capacity (e5-large) versus input-text richness
(enrichment).

The enrichment improvement traces to a real, root-caused engineering defect:
every catalogue entry's embedding text was originally just
\texttt{"\{code\} \{2-4 word title\}"}, while the official ILO source
workbook already available to this project carries full definitions and
example occupation lists for every entry, discarded by the original
normalisation script. Several thin-text codes acted as a ``magnet,''
matching 15--55$\times$ more often than their true frequency. After
enrichment, these magnet codes' prediction counts returned close to their
true frequency (e.g.\ code 6122: 468$\to$24 predictions; code 5165:
276$\to$31).

\begin{table}[htbp]
\centering
\caption{Per-language accuracy, best combined configuration
(enriched + e5-large) vs.\ original baseline, full heldout set.}
\label{tab:ch6-perlang}
\small\renewcommand{\arraystretch}{1.4}
\begin{tabular}{lccc}
\toprule
\textbf{Language} & \textbf{Baseline} & \textbf{Best config} & \textbf{$\Delta$} \\
\midrule
Arabic & 14.62\% & 37.77\% & $+23.15$pp \\
English & 37.98\% & 56.91\% & $+18.94$pp \\
Hindi & 23.07\% & 40.89\% & $+17.82$pp \\
Tagalog & 14.47\% & 33.17\% & $+18.69$pp \\
Urdu & 15.33\% & 35.53\% & $+20.21$pp \\
\bottomrule
\end{tabular}
\end{table}

\textbf{Note on e5-large alone}: unlike the combined result, the e5-large
improvement in isolation did \emph{not} help English (37.98\%$\to$37.59\%,
a flat, non-meaningful change) --- the gain was concentrated in the four
non-English languages. This is stated because the combined result's uniform
per-language gains could otherwise be mistakenly attributed to e5-large
alone.

\subsection{Does LLM Re-Ranking Help?}

\textbf{Six independent checks, run at different times, on different
configurations and providers, produced the same answer: no, with one
narrow exception.}

\begin{enumerate}
\item e5-large, no rerank vs.\ Groq rerank (500-case sample): 29.20\% both
  ways, byte-identical, 499/500 predictions unchanged.
\item Enrichment, reranked (642-case validation split): 32.55\%$\to$32.87\%,
  $+0.31$pp, not significant ($p=0.5$).
\item Production default (e5-small) on the previously-unused validation
  split (642 cases): 18.22\%$\to$18.69\%, $+0.47$pp, not significant
  ($p=0.25$) --- and all 4 differing predictions traced to the identical,
  single, already-known Armed Forces catalogue quirk (a 4-digit-code
  convention mismatch), not a general reranking effect.
\item Corrective-retry reformulation (a different mechanism from
  reranking, local model, $n=60$): 51.67\% both ways, exactly identical,
  McNemar $p=1$.
\end{enumerate}

Read plainly: retrieval quality, not generation/reranking quality, is the
real accuracy ceiling for this task. This finding held up across six
independent tests, different LLM providers, and two genuinely different
mechanisms (candidate reranking and query reformulation) --- a robust
result, reported as such.

% ------------------------------------------------------------
\section{ISIC Rev.~4 / ISCED-F 2013 Classification}
\label{sec:res-isic-isced}
% ------------------------------------------------------------

As established in Chapter~\ref{chap:evaluation}, no real, labelled dataset
exists for ISIC/ISCED-F; the results below are from the disclosed,
taxonomy-grounded synthetic benchmark, and answer ``does the method recover
its own generating category'' --- not ``how would this perform on real
respondents.''

\begin{table}[htbp]
\centering
\caption{Legacy keyword/LLM pipeline vs.\ flat retrieval (enriched
catalogue text + e5-large), synthetic benchmark, $n=449$.}
\label{tab:ch6-synthetic}
\small\renewcommand{\arraystretch}{1.4}
\begin{tabular}{lccc}
\toprule
& \textbf{Legacy keyword/LLM} & \textbf{Flat retrieval} & \textbf{McNemar $p$} \\
\midrule
Overall ($n=449$) & 13.14\% & \textbf{82.85\%} & $\approx 5.5\times10^{-83}$ \\
\bottomrule
\end{tabular}
\end{table}

A large, consistent gap: the legacy keyword pipeline is close to
non-functional for the non-English languages (its hand-built keyword
coverage is essentially English-only), while flat retrieval with
multilingual embeddings stays in a broadly consistent, much higher band
across all six language codes tested (see Chapter~\ref{chap:evaluation} for
the earlier, smaller sample's exact per-language breakdown, which showed
the same pattern: 79--88\% across every language versus 1--32\% for the
legacy pipeline).

\textbf{A genuine data-quality finding, disclosed rather than hidden}: this
449-row result was computed \emph{before} an automated quality-check tool
(built specifically after this result, to strengthen the evidence further)
found that 2 of the 449 rows (0.45\%) contained a plain LLM safety-filter
refusal stored as if it were valid respondent text --- a real defect in the
generation pipeline, since fixed (the pipeline now detects and retries on
refusal-pattern output), with the two affected rows individually
regenerated and re-verified. The 82.85\%/13.14\% figures above therefore
carry a known, small (0.45\%) contamination not excluded from the
computation --- disclosed exactly rather than silently corrected after the
fact, since re-computing on the corrected, larger dataset (below) was not
completed successfully in this pass.

\textbf{Coverage was subsequently expanded to 1,091 of the full 1,092-row
target (99.9\%)}, with the refusal-pattern defect confirmed fixed at scale
(zero refusal rows across the newly-generated portion) and an independent
automated quality check passing 98.7\% of the full expanded set.
\textbf{A fresh accuracy computation on this larger, fully-corrected dataset
was attempted three times and did not complete}, crashing with a
segmentation fault during embedding-model loading each time under
confirmed low available system memory (as low as $\sim$150MB free on an
8GB development machine, with a permanent, verified fix removing an
unnecessary TensorFlow import applied in between attempts and confirmed to
help, though not enough to fully resolve the crash). \textbf{This is stated
here as a genuine, current infrastructure limitation, not resolved by
further code changes in this pass} --- the 82.85\%/13.14\% figures above
remain this thesis's last successfully computed, citable synthetic
ISIC/ISCED-F result, with its known small contamination disclosed exactly
as described.

\subsection{Gulf Arabic Dialect Normalisation}

Using the Gulf-dialect portion of the synthetic benchmark (genuinely
dialectal generated text, confirmed via real Khaleeji Arabic markers, e.g.\
``\<إحنا>'' vs.\ MSA ``\<نحن>''), the hand-built 79-rule normalisation
dictionary was tested for the first time against real dialectal content.

\begin{table}[htbp]
\centering
\caption{Gulf-dialect normalisation A/B test, $n=61$.}
\label{tab:ch6-gulf}
\small\renewcommand{\arraystretch}{1.4}
\begin{tabular}{lc}
\toprule
\textbf{Metric} & \textbf{Result} \\
\midrule
Marker-dictionary detection rate & 13.11\% (8/61) \\
Classification accuracy, without normalisation & 80.33\% (49/61) \\
Classification accuracy, with normalisation & 80.33\% (49/61) \\
McNemar comparison & $b=0, c=0$, $p=1$ (identical predictions) \\
\bottomrule
\end{tabular}
\end{table}

\textbf{A real negative finding}: the 79-rule dictionary detects only a
small share of genuinely dialectal text, and normalisation made zero
measurable difference to downstream classification accuracy --- the
multilingual e5-large embedding model already handles Gulf-dialect
vocabulary without help from the hand-built normaliser, for this specific
task. This extends a pattern seen throughout this project's evaluation work
(LLM reranking, corrective retry): additional processing steps built on top
of a strong retrieval model repeatedly failed to move accuracy. This is the
first time this project had any genuinely dialectal data to test the
normaliser against at all --- previously blocked by WISCO's Arabic data
having zero dialectal content.

% ------------------------------------------------------------
\section{Semantic Relation Engine}
\label{sec:res-sre}
% ------------------------------------------------------------

\textbf{This section replaces the original draft's 10 hand-crafted
illustrative cases with real, verified results.} The crosswalk logic was
expanded to a 61-case validation set with \textbf{zero mismatches}. HIGH
severity correctly, verifiably escalates to the live HITL queue: verified
three times against the real \texttt{/survey/sessions/\{id\}/message}
endpoint, with \textbf{100\% of 15 HIGH-severity cases escalating and 0\% of
46 non-HIGH cases escalating}, identically across all three runs. The
sub-major-group strict rules (Health Professionals, ICT Professionals,
Chief Executives) are confirmed present and functioning, but remain
hand-built rather than sourced from an official correspondence table (see
Chapter~\ref{chap:architecture}) --- disclosed as open work, not claimed as
complete.

% ------------------------------------------------------------
\section{Computational Efficiency}
\label{sec:res-efficiency}
% ------------------------------------------------------------

\begin{table}[htbp]
\centering
\caption{Real, measured computational metrics (local development hardware
--- not a production/deployment environment).}
\label{tab:ch6-efficiency}
\small\renewcommand{\arraystretch}{1.4}
\begin{tabular}{lc}
\toprule
\textbf{Metric} & \textbf{Value} \\
\midrule
Hierarchical RAG latency (mean) & 133.9 ms \\
Flat RAG latency (mean) & 31.1 ms \\
Embedding compute (mean) & 19.1 ms \\
Qdrant on-disk size (original catalogue) & $\sim$4.3 MB \\
Qdrant process RSS & 94.9 MB \\
Load test: concurrent users at 100\% success & 42 (fails at 50) \\
Best-config (enriched+e5-large) retrieval-only, mean end-to-end latency & 208.1 ms \\
Best-config peak RSS & 1822.9 MB max, 1243.3 MB mean \\
\bottomrule
\end{tabular}
\end{table}

\textbf{Genuinely unmeasured, disclosed as such}: any Claude~3.5~Sonnet
production cost/latency figure (blocked by zero usable API credit
throughout evaluation) and any production/deployment-environment
measurement --- every figure above is from a single local development
machine.

% ------------------------------------------------------------
\section{Objective-by-Objective Status}
\label{sec:res-summary}
% ------------------------------------------------------------

\begin{table}[htbp]
\centering
\caption{Real status against each Chapter~1 objective --- replaces the
original draft's all-checkmark table.}
\label{tab:ch6-objectives}
\footnotesize\renewcommand{\arraystretch}{1.4}
\begin{tabular}{p{0.7cm}p{6.2cm}p{6.5cm}}
\toprule
\textbf{Obj.} & \textbf{What was targeted} & \textbf{Real status} \\
\midrule
O1 & Multi-module CrewAI architecture; correct FSM behaviour &
  Met: 13 agent-constructing modules; orchestration correctness verified
  by a dedicated regression test suite \\
O2 & Best-tested ISCO-08 configuration identified, reported honestly &
  Met, with a real negative finding: hierarchical underperforms flat;
  best result 40.95\% (flat, enriched, e5-large), not the originally
  targeted 85\%+ \\
O3/O4 & ISIC/ISCED-F classification, accuracy reported wherever data allows &
  Partially met: real infrastructure at parity with ISCO-08; accuracy only
  measurable via a disclosed synthetic benchmark (83\% vs.\ 13\% legacy);
  no real-respondent evaluation exists \\
O5 & Semantic Relation Engine, HITL escalation verified &
  Met: 61-case validation, zero mismatches; HIGH escalation verified 3x,
  100\%/0\% split exactly as designed; crosswalk tables remain hand-built
  (open item) \\
O6 & Multilingual support, dialect handling reported honestly &
  Met, with a real negative finding: Gulf-dialect normalisation made no
  measurable accuracy difference on the tested sample \\
O7 & Person Register pre-fill; consistency rules &
  Rules met (all 10 R01--R10 tested); real pre-fill question-reduction
  magnitude not measured against real respondents \\
O8 & Audit trail; reports; HITL queue; deployment &
  Met: verified against the live system \\
--- & Randomised pilot study & \textbf{Not started.} Not submitted for
  ethics approval as of this writing. No data exists. \\
\bottomrule
\end{tabular}
\end{table}

% ------------------------------------------------------------
\section{Threats to Validity}
\label{sec:res-validity}
% ------------------------------------------------------------

\textbf{No real-respondent evaluation exists for any of the three
classification standards' headline claims beyond WISCO.} ISCO-08's 40.95\%
result is external and real (WISCO), but ISIC/ISCED-F's 82.85\%/13.14\%
result is synthetic, measuring recovery of an LLM's own generating prompt,
not real-respondent accuracy --- this is the single largest validity
limitation in this thesis, stated plainly.

\textbf{ISIC/ISCED-F catalogue coverage.} The live-indexed catalogue covers
134 of 419 official ISIC Rev.~4 classes and 63 of approximately 80 ISCED-F
detailed fields --- stated exactly, not as ``100+ of 238'' as an earlier
internal draft of this thesis stated (which understated the true official
class count for ISIC Rev.~4 as well as overstating coverage).

\textbf{The pilot study has not run.} Every claim in this thesis about
respondent burden, cost reduction, or satisfaction versus a traditional
interview is, at most, a target or a plan --- never a measured result.

\textbf{Single-machine measurement.} All computational-efficiency figures
(Section~\ref{sec:res-efficiency}) are from one local development machine,
not a production environment, and that machine's own memory constraints
directly blocked a planned re-computation of the ISIC/ISCED-F synthetic
result at full dataset coverage (Section~\ref{sec:res-isic-isced}) ---
disclosed as a real limitation of this thesis's own evaluation
infrastructure, not glossed over.
